\section{问题与命题}

\begin{frame}{海缆自主巡检的核心矛盾}
  \begin{columns}[T,onlytextwidth]
    \begin{column}{0.48\textwidth}
      \begin{block}{环境与观测}
        \begin{itemize}
          \item DVL 丢包、声呐杂波、磁畸变随工况变化
          \item 多源传感器异步、频率相差一个数量级
          \item 掩埋电缆不可见，单一模态存在失效区
        \end{itemize}
      \end{block}
    \end{column}
    \begin{column}{0.48\textwidth}
      \begin{alertblock}{传统链路缺口}
        \begin{itemize}
          \item 状态估计只输出“在哪里”
          \item 决策与控制不知道“有多可信”
          \item 观测退化时仍可能维持原有控制强度
        \end{itemize}
      \end{alertblock}
    \end{column}
  \end{columns}

  \FrameConclusion{矛盾不只是“估得准不准”，而是系统能否在估计变差时主动收敛行为。}
\end{frame}

\begin{frame}{本文的可检验命题}
  \centering
  \vspace{3mm}
  {\Large
    观测质量
    \(\longrightarrow\)
    协方差
    \(\longrightarrow\)
    置信度
    \(\longrightarrow\)
    决策阈值与控制权重
  }

  \vspace{8mm}
  \begin{columns}[T,onlytextwidth]
    \begin{column}{0.31\textwidth}
      \begin{block}{H1\quad 可观测}
        扰动增强时，自适应协方差应被触发且可核查。
      \end{block}
    \end{column}
    \begin{column}{0.31\textwidth}
      \begin{block}{H2\quad 可调度}
        低置信度应改变行为优先级和控制代价。
      \end{block}
    \end{column}
    \begin{column}{0.31\textwidth}
      \begin{block}{H3\quad 有边界}
        感知完全退化时，权重调整不应被宣称仍有优势。
      \end{block}
    \end{column}
  \end{columns}

  \FrameConclusion{验证目标是“机制、收益与失效边界”同时成立，而非追求单一最优数字。}
\end{frame}
